\begin{theorem}\label{thm:complex-npp-ledger}
Let \(1\leq p\leq n\).  Whenever the local mass condition in
\Cref{sec:ch51-npp-local} holds, the product in
\eqref{eq:complex-npp-truncation} is a closed positive \((p,p)\)-current.
It is symmetric, local for the plurifine topology, depends only on the
currents, agrees with the Bedford--Taylor product for locally bounded
potentials, and charges no pluripolar set.  It is multilinear on the positive
cone in the following precise sense: if \(S_1,T_1,T_2,\ldots,T_p\) are
closed positive \((1,1)\)-currents, then
\[
 \left\langle (S_1+T_1)\wedge T_2\wedge\cdots\wedge T_p\right\rangle
 =\left\langle S_1\wedge T_2\wedge\cdots\wedge T_p\right\rangle
  +\left\langle T_1\wedge T_2\wedge\cdots\wedge T_p\right\rangle,
\]
where the left-hand side is locally finite if and only if both terms on the
right are locally finite.

On a compact K\"ahler manifold the local mass condition holds for every
tuple of global closed positive \((1,1)\)-currents, so
\eqref{eq:complex-global-npp} is always defined.  If
\(T_j=\theta_j+\ddc\varphi_j\) and
\(T_j'=\theta_j+\ddc\psi_j\), with
\(\varphi_j\preceq\psi_j\) for all \(j\), then, for every K\"ahler form
\(\omega\),
\[
 \int_X\left\langle T_1\wedge\cdots\wedge T_p\right\rangle
              \wedge\omega^{n-p}
 \leq
 \int_X\left\langle T_1'\wedge\cdots\wedge T_p'\right\rangle
              \wedge\omega^{n-p}.
\]
More generally, suppose that \(X\) is a compact complex manifold admitting
a modification \(\mu:Y\to X\) from a compact K\"ahler manifold \(Y\).
Then the local mass condition holds for every tuple of closed positive
\((1,1)\)-currents \(T_1,\ldots,T_p\) on \(X\), their pullbacks are
well defined as closed positive currents, and
\[
 \left\langle T_1\wedge\cdots\wedge T_p\right\rangle
 =\mu_*\left\langle
       \mu^*T_1\wedge\cdots\wedge\mu^*T_p\right\rangle .
\tag{51.16a}\label{eq:complex-npp-modification}
\]
In particular, for \(p=n\),
\[
 \int_X\left\langle T_1\wedge\cdots\wedge T_n\right\rangle
 =\int_Y\left\langle
       \mu^*T_1\wedge\cdots\wedge\mu^*T_n\right\rangle .
\]
\end{theorem}

\begin{proof}
\emph{Positivity, compatibility, and the bounded case.}
The restricted currents in \eqref{eq:complex-npp-truncation} are positive
and increase with \(k\).  Under the local mass condition their limit is
therefore a positive current.  Plurifine locality makes the truncations
compatible and also shows that two tuples agreeing on a plurifine open set
have the same product there.  If all potentials are locally bounded, then
on each relatively compact set \(O_k\) is the whole set for all sufficiently
large \(k\), and the construction reduces to the Bedford--Taylor product.

\emph{No mass on pluripolar sets.}  Assume that the increasing limit in
\eqref{eq:complex-npp-truncation} is locally finite.  In the compact K\"ahler
situation of \eqref{eq:complex-global-npp}, this local finiteness follows from
\cite[Proposition~1.6]{BEGZ10}.  Let \(A\) be a Borel pluripolar set.  Work
on coordinate balls \(V\Subset V'\Subset U\), with \(V'\) strongly
pseudoconvex.  Choose a smooth strictly psh function \(\rho\) on a
neighbourhood of \(\overline{V'}\), and set \(\beta:=\ddc\rho\).  Recall
that the Bedford--Taylor capacity and its outer capacity are
\[
 \Capa_{\mathrm{BT}}(E,V')
 :=\sup_{\substack{q\in\PSH(V')\\0<q<1}}
       \int_E(\ddc q)^n,
 \qquad
 \Capa_{\mathrm{BT}}^*(E,V')
 :=\inf_{\substack{G\supset E\\G\ \mathrm{open}}}
       \Capa_{\mathrm{BT}}(G,V').
\]
For fixed \(k\), every
\(w_{j,k}:=\max\{u_j,-k\}\) is locally bounded and psh.  The mixed
Bedford--Taylor measure
\[
 \nu_k:=\one_{O_k}\bigwedge_{j=1}^p\ddc w_{j,k}
                   \wedge\beta^{n-p}
\]
vanishes on \(A\cap V\).  Indeed, each \(w_{j,k}\) is bounded on \(V'\)
because \(\overline{V'}\Subset U\), and \(\rho\) is bounded there because
it is smooth on a neighbourhood of \(\overline{V'}\).  After adding
constants and multiplying by positive scalars, normalize the \(p\) functions
\(w_{j,k}\) and the \(n-p\) copies of \(\rho\) to functions with values in
\((0,1)\) on \(V'\).  Denote the normalized functions by
\(f_i=a_i h_i+b_i\), where
\((h_1,\ldots,h_n)=(w_{1,k},\ldots,w_{p,k},\rho,\ldots,\rho)\) and
\(a_i>0\), and put \(W=n^{-1}\sum_i f_i\).  Positivity and
multilinearity give
\[
 (\ddc W)^n\geq\frac{n!}{n^n}\bigwedge_{i=1}^n\ddc f_i
 =\frac{n!}{n^n}\left(\prod_{i=1}^na_i\right)
   \bigwedge_{j=1}^p\ddc w_{j,k}\wedge\beta^{n-p}.
\]
Consequently, for Borel \(E\subset V\),
\[
 \nu_k(E)\leq C_k\int_E(\ddc W)^n
             \leq C_k\Capa_{\mathrm{BT}}^*(E,V').
\]
Here the first inequality also uses the fact that restriction by
\(\one_{O_k}\) can only decrease a positive measure.  A pluripolar set has
zero Bedford--Taylor outer capacity
\cite[Theorem~6.9]{BT82}.  Finally, the currents in
\eqref{eq:complex-npp-truncation} increase to the nonpluripolar product, so
their trace measures \(\nu_k\) increase to its trace measure on \(V\).
Continuity from below therefore gives
\[
 \left(\left\langle\bigwedge_{j=1}^p\ddc u_j\right\rangle
       \wedge\beta^{n-p}\right)(A\cap V)
 =\lim_k\nu_k(A\cap V)=0.
\]
Since \(V\) is arbitrary, the product charges no pluripolar set.

\emph{Independence of potentials and multilinearity.}
The two mass comparisons
\eqref{eq:complex-npp-potential-mass-forward} and
\eqref{eq:complex-npp-potential-mass-reverse} prove that local finiteness is
unchanged when \(u_j\) is replaced by \(u'_j=u_j+h_j\) with \(h_j\)
pluriharmonic.  The capacity argument above has now established the
no-charge property without using potential independence.  Hence the precise
admissible-exhaustion rule
\eqref{eq:complex-npp-admissible-exhaustion}, applied to the two sides of
\eqref{eq:complex-npp-potential-common-current}, proves equality of the two
limits on every relatively compact common coordinate open set.  This proves
both potential independence and the patching asserted in
\eqref{eq:complex-global-npp}.

We give the corresponding argument for addition, since truncation itself is
not linear.  Write locally
\(S_1=\ddc u\), \(T_1=\ddc v\), and
\(T_j=\ddc u_j\) for \(j\geq2\); after subtracting constants, suppose that
all these psh functions are nonpositive.  Put
\[
 E_k=\{u>-k\}\cap\{v>-k\}
             \cap\bigcap_{j=2}^p\{u_j>-k\}.
\]
On \(E_k\), plurifine locality and Bedford--Taylor multilinearity give
\[
\begin{aligned}
 &\one_{E_k}\ddc\max\{u+v,-2k\}
       \wedge\bigwedge_{j=2}^p\ddc\max\{u_j,-2k\}\\
 &\quad=
 \one_{E_k}\ddc\max\{u,-k\}
       \wedge\bigwedge_{j=2}^p\ddc\max\{u_j,-k\}\\
 &\qquad\quad+
 \one_{E_k}\ddc\max\{v,-k\}
       \wedge\bigwedge_{j=2}^p\ddc\max\{u_j,-k\}.
\end{aligned}
\tag{51.16b}\label{eq:complex-npp-common-exhaustion}
\]
The sets \(E_k\) increase, and their union has pluripolar complement.  Once
the three products are locally finite, the admissible-exhaustion rule and
\eqref{eq:complex-npp-common-exhaustion} give the desired equality.

It remains to justify the asserted equivalence of local finiteness.  If the
two products on the right are locally finite, then the canonical level-
\(m\) locus for \(u+v,u_2,\ldots,u_p\) is contained in \(E_m\), because
\(u,v\leq0\) and \(u+v>-m\) imply both \(u>-m\) and \(v>-m\).
Equation \eqref{eq:complex-npp-common-exhaustion} and positivity therefore
bound its mass by the sum of the two right-hand masses.  Conversely, suppose
the product on the left is locally finite.  Fix \(m\).  The level-\(m\)
truncation for \(u,u_2,\ldots,u_p\) charges no pluripolar set, and hence its
mass on a compact set is the limit of its masses after intersection with
\(\{v>-k\}\).  On that intersection, for \(k\geq m\), positivity and
Bedford--Taylor multilinearity bound it by the level-\((m+k)\) truncation
for \(u+v,u_2,\ldots,u_p\).  The latter masses are uniformly bounded by
hypothesis.  Letting \(k\to\infty\), and then taking the supremum over
\(m\), proves local finiteness of the product with \(u\); the argument with
\(v\) is identical.  This proves the full assertion of
\cite[Proposition~1.4]{BEGZ10} without claiming a false levelwise
truncation identity.

\emph{Global local finiteness.}
Let \(\omega\) be a K\"ahler form on the compact K\"ahler manifold \(X\).
By the just-proved additivity, adding a sufficiently large multiple of
\(\omega\) to each \(T_j\) reduces the local-finiteness assertion to the
case in which every class \(\{T_j\}\) is K\"ahler.  Choose a K\"ahler form
\(\omega_j\in\{T_j\}\) and
\(\varphi_j\in\PSH(X,\omega_j)\), normalized by \(\varphi_j\leq0\), such
that \(T_j=\omega_j+\ddc\varphi_j\).  On a sufficiently small coordinate
open set \(U\), write \(\omega_j=\ddc\psi_j\), with \(\psi_j\leq0\) smooth
and psh, and put \(u_j=\psi_j+\varphi_j\).  On
\(O_k=\bigcap_j\{u_j>-k\}\) one has \(\varphi_j>-k\), and the two bounded
psh functions
\[
       \max\{u_j,-k\}
       \quad\hbox{and}\quad
       \psi_j+\max\{\varphi_j,-k\}
\]
coincide.  Therefore, for every compact \(K\Subset U\),
\[
\begin{aligned}
 &\int_{K\cap O_k}\omega^{n-p}\wedge
       \bigwedge_{j=1}^p\ddc\max\{u_j,-k\}\\
 &\quad\leq
 \int_X\omega^{n-p}\wedge
       \bigwedge_{j=1}^p
          \bigl(\omega_j+\ddc\max\{\varphi_j,-k\}\bigr)\\
 &\quad=
 \int_X\omega^{n-p}\wedge\omega_1\wedge\cdots\wedge\omega_p.
\end{aligned}
\tag{51.16c}\label{eq:complex-npp-global-finiteness}
\]
The inequality uses positivity; the last equality is Stokes' theorem for
the globally bounded \(\omega_j\)-psh truncations.  The bound is independent
of \(k\), proving the local mass condition.  This is the argument of
\cite[Proposition~1.6]{BEGZ10}; in particular it does not use the generally
false positivity of
\(\theta_j+\ddc\max\{\varphi_j,-k\}\) for an arbitrary smooth
representative \(\theta_j\).

\emph{Closedness.}
We reproduce the local argument of \cite[Theorem~1.8]{BEGZ10}, since the
restricted truncations are not closed term by term.  The assertion is
automatic for \(p=n\).  For \(p<n\), work near the origin in
\(\mathbf C^n\), write \(T_j=\ddc u_j\) with \(u_j\leq0\), and let
\(\rho\) be a positive constant-coefficient
\((n-p-1,n-p-1)\)-form.  Set
\[
 \Theta_k:=\rho\wedge\bigwedge_{j=1}^p
                  \ddc\max\{u_j,-k\},
 \qquad O_k:=\bigcap_{j=1}^p\{u_j>-k\}.
\]
Then \(\Theta_k\) is a closed positive current of bidimension \((1,1)\),
and \(\one_{O_k}\Theta_k\) converges to
\(\rho\wedge\langle T_1\wedge\cdots\wedge T_p\rangle\).  Put
\(u=\sum_j u_j\).  Choose a smooth convex nondecreasing function \(\chi\)
with \(\chi(t)=0\) for \(t\leq1/2\) and \(\chi(1)=1\), and define
\[
                         w_k=\chi(\rme^{u/k}).
\]
Thus \(0\leq w_k\leq1\), the bounded psh functions \(w_k\) increase to
\(1\) off \(\{u=-\infty\}\), and \(w_k=0\) outside \(O_k\), since
\(\rme^{-1}<1/2\).  Choose a smooth nondecreasing function
\(0\leq\vartheta\leq1\) that vanishes on \((-\infty,1/2]\) and equals
\(1\) near \(1\).  Then
\(\vartheta(w_k)\leq\one_{O_k}\), while
\(\vartheta'(w_k)\) vanishes outside \(O_k\) and converges to zero off
\(\{u=-\infty\}\).

By plurifine locality,
\(\one_{O_k}\Theta_k\) is the restriction to \(O_k\) of the limiting
current.  Hence
\[
 0\leq(\one_{O_k}-\vartheta(w_k))\Theta_k
 \leq(1-\vartheta(w_k))
       \rho\wedge\left\langle T_1\wedge\cdots\wedge T_p\right\rangle .
\]
The right-hand side tends weakly to zero by dominated convergence and the
no-charge property already proved.  It is therefore enough to show that
\(\mathrm d(\vartheta(w_k)\Theta_k)\to0\).

For a bounded psh function \(w\), a smooth function \(f\), and a closed
positive current \(\Theta\) of bidimension \((1,1)\), the Bedford--Taylor
chain rule is
\[
             \mathrm d(f(w)\Theta)=f'(w)\,\mathrm d w\wedge\Theta.
\tag{51.16d}\label{eq:complex-npp-chain-rule}
\]
To verify it, decrease smooth psh functions \(w_\ell\) to \(w\).  The
Bedford--Taylor monotone convergence theorem gives convergence of
\(w_\ell\Theta\) and of the energy currents
\(\mathrm d w_\ell\wedge\mathrm d^c w_\ell\wedge\Theta\).  For a test
one-form \(\eta\), write
\[
 \bigl(f'(w_\ell)-f'(w)\bigr)\,\mathrm d w_\ell\wedge\Theta
 +f'(w)\,\mathrm d(w_\ell-w)\wedge\Theta
\]
for the difference of the derivative terms.  Cauchy--Schwarz bounds the
first summand by the product of a uniformly bounded energy integral and an
integral containing \(\lvert f'(w_\ell)-f'(w)\rvert^2\), which tends to
zero by dominated convergence.  It bounds the second by a fixed factor
times
\[
 \int\tau\,\mathrm d(w_\ell-w)\wedge
                    \mathrm d^c(w_\ell-w)\wedge\Theta,
\]
which tends to zero by Bedford--Taylor convergence; here \(\tau\) is a
cutoff equal to \(1\) on the support of \(\eta\).  Passing to the limit in
the smooth chain rules proves \eqref{eq:complex-npp-chain-rule}.

Apply this with \(w=w_k\), \(f=\vartheta\), and \(\Theta=\Theta_k\).
For a compactly supported smooth one-form \(\eta\), choose a nonnegative
cutoff \(\tau\) equal to \(1\) on its support.  Cauchy--Schwarz for the
positive current \(\Theta_k\) bounds the square of the pairing with
\(\vartheta'(w_k)\,\mathrm d w_k\wedge\Theta_k\) by a fixed multiple of
the product
\[
 \left(\int\tau\,\mathrm d w_k\wedge\mathrm d^c w_k\wedge\Theta_k\right)
 \left(\int\vartheta'(w_k)^2\one_{O_k}\Theta_k\wedge\omega\right),
\]
where \(\omega\) is the Euclidean K\"ahler form.  The first factor is
uniformly bounded: positivity gives
\[
 2\int\tau\,\mathrm d w_k\wedge\mathrm d^c w_k\wedge\Theta_k
 \leq\int\tau\,\ddc(w_k^2)\wedge\Theta_k
 =\int_{O_k}w_k^2\,\ddc\tau\wedge\Theta_k,
\]
and the absolute value of the last integral is controlled by the locally
finite mass of \(\one_{O_k}\Theta_k\) on a fixed compact set.  The second
factor tends to zero, since
\[
 \vartheta'(w_k)^2\one_{O_k}\Theta_k
 \leq \vartheta'(w_k)^2
       \rho\wedge\left\langle T_1\wedge\cdots\wedge T_p\right\rangle
\]
and dominated convergence applies.  Thus
\(\mathrm d(\vartheta(w_k)\Theta_k)\to0\).  Constant-coefficient forms
\(\rho\) span the required test directions, so the nonpluripolar product is
closed.

\emph{Monotonicity of total masses.}
The displayed inequality in the statement is not a pointwise inequality of
currents, nor does it follow from integration by parts on the restricted
truncations.  In the stated generality it is exactly
\cite[Theorem~1.1]{DDNL18mono}.  The earlier result
\cite[Theorem~1.16]{BEGZ10} assumes small unbounded locus.  This chapter
imports the unrestricted theorem as
item \cref{itm:ch51-import:monotone}.  When \(p<n\), its \(n\)-fold mixed
statement is applied after adjoining \(n-p\) copies of the smooth K\"ahler
current \(\omega\), with zero potential.

\emph{Modifications.}
We first construct the pullback occurring in the statement.  On a coordinate
open set \(U\subset X\), choose a psh potential \(u_j\) such that
\(T_j|_U=\ddc u_j\), and set
\[
 \left.\mu^*T_j\right|_{\mu^{-1}(U)}
       :=\ddc(u_j\circ\mu).
\tag{51.16e}\label{eq:complex-npp-current-pullback}
\]
A modification is biholomorphic over a dense open set, and every component of
\(Y\) meets its inverse image.  Thus \(u_j\circ\mu\) is not identically
\(-\infty\) on any component, and is therefore psh.  On an overlap, two
choices of \(u_j\) differ by a
pluriharmonic function; their compositions with the holomorphic map
\(\mu\) are again pluriharmonic.  Thus
\eqref{eq:complex-npp-current-pullback} patches and defines a closed positive
\((1,1)\)-current on \(Y\).

Choose the standard analytic exceptional pair: a nowhere-dense closed
analytic set \(Z\subset X\) and
\[
                         F:=\mu^{-1}(Z)\subset Y
\]
such that \(\mu:Y\setminus F\to X\setminus Z\) is biholomorphic.  Since
\(Y\) is compact K\"ahler, the already proved global finiteness gives the
closed positive current
\[
 R_Y:=\left\langle
       \mu^*T_1\wedge\cdots\wedge\mu^*T_p\right\rangle,
 \qquad S:=\mu_*R_Y .
\]
Both \(F\) and \(Z\) are pluripolar, and the no-charge property gives zero
trace mass for \(R_Y\) on \(F\).  Here is the corresponding assertion for
the pushforward, including the trace forms.  Choose a smooth Hermitian form
\(\sigma\) on \(X\) and a K\"ahler form \(\eta\) on \(Y\).  Compactness gives
\(\mu^*\sigma\leq A\eta\) for some \(A>0\), and positivity gives
\[
 0\leq R_Y\wedge(\mu^*\sigma)^{n-p}
 \leq A^{n-p}R_Y\wedge\eta^{n-p}.
\]
The projection formula for the smooth form \(\sigma^{n-p}\) consequently
shows
\[
\begin{aligned}
 (S\wedge\sigma^{n-p})(Z)
 &=\bigl(\mu_*(R_Y\wedge(\mu^*\sigma)^{n-p})\bigr)(Z)\\
 &=\bigl(R_Y\wedge(\mu^*\sigma)^{n-p}\bigr)(F)=0.
\end{aligned}
\tag{51.16f}\label{eq:complex-npp-pushforward-zero-trace}
\]

It remains to prove, rather than assume, the local mass condition downstairs.
On a coordinate open set \(U\subset X\), use the potentials \(u_j\) above,
put \(O_k=\bigcap_j\{u_j>-k\}\), and denote the \(k\)-th restricted bounded
product by
\[
 R_{X,k}:=\one_{O_k}\bigwedge_{j=1}^p
                         \ddc\max\{u_j,-k\}.
\]
On \(U\setminus Z\), biholomorphic invariance and plurifine locality identify
\(R_{X,k}\) with the corresponding stage upstairs; in particular
\(R_{X,k}\leq S\) there.  The fixed-level part of the Bedford--Taylor capacity
argument above, which uses only bounded psh functions and not local finiteness
of the limiting product, shows that \(R_{X,k}\) has zero trace mass on the
analytic set \(Z\): on each relatively compact coordinate ball, the Hermitian
trace form \(\sigma\) is uniformly comparable with the strictly psh trace
form used in that argument.  Therefore, for every compact \(K\Subset U\),
\[
 \int_{K\cap O_k}\sigma^{n-p}\wedge
       \bigwedge_{j=1}^p\ddc\max\{u_j,-k\}
 =\int_{K\setminus Z}R_{X,k}\wedge\sigma^{n-p}
 \leq\int_K S\wedge\sigma^{n-p}.
\tag{51.16g}\label{eq:complex-npp-modification-local-mass}
\]
The right-hand side is finite and independent of \(k\).  Hence the local
mass condition holds on \(X\), and the increasing limit \(R_X\) of the
\(R_{X,k}\) is a locally finite current.  On \(X\setminus Z\), the preceding
biholomorphic identification gives \(R_X=S\).  The no-charge property for
\(R_X\), together with \eqref{eq:complex-npp-pushforward-zero-trace}, says
that neither current has trace mass on \(Z\); equality therefore extends
across \(Z\).  This proves
\(R_X=S\), which is exactly \eqref{eq:complex-npp-modification} and the
formula of \cite[Remark~1.7]{BEGZ10}.

The top-degree mass equality follows by integrating.  More generally, for
any smooth \((n-p,n-p)\)-form \(\gamma\), equality of currents and the
projection formula give
\[
 \int_X\left\langle T_1\wedge\cdots\wedge T_p\right\rangle\wedge\gamma
 =\int_Y\left\langle\mu^*T_1\wedge\cdots\wedge\mu^*T_p\right\rangle
              \wedge\mu^*\gamma .
\]
\end{proof}
